\chapter*{LỜI CẢM ƠN}
\addcontentsline{toc}{chapter}{LỜI CẢM ƠN}

Lời đầu tiên, nhóm em xin gửi lời cảm ơn chân thành và sâu sắc nhất đến ThS. Đỗ Quang Thơ – giảng viên trực tiếp giảng dạy và hướng dẫn học phần Thiết kế \& Kiến trúc phần mềm. Trong suốt quá trình học tập và thực hiện bài tập lớn, thầy đã tận tình truyền đạt những kiến thức nền tảng về kiến trúc phần mềm, từ các nguyên lý thiết kế cơ bản như SOLID, mô hình phân tầng cho đến các kiến trúc hiện đại như Microservices và Event- Driven Architecture. Những kiến thức này không chỉ giúp nhóm hoàn thành bài tập lớn mà còn là hành trang quý báu cho sự nghiệp phát triển phần mềm sau này của mỗi thành viên. Nhóm em cũng xin gửi lời cảm ơn đến Khoa Công nghệ Thông tin – Trường Đại học Đại Nam đã tạo điều kiện thuận lợi về cơ sở vật chất và môi trường học tập, giúp nhóm có thể tập trung nghiên cứu và phát triển dự án một cách hiệu quả nhất. Hệ thống thư viện, phòng máy và các tài nguyên học thuật của trường là những hỗ trợ không thể thiếu trong suốt quá trình thực hiện. Đặc biệt, nhóm xin cảm ơn sự phối hợp nhịp nhàng và hiệu quả từ nhóm N1 (Catalog Service) và nhóm N3 (Identity \& Report Service). Việc xây dựng thành công một hệ thống microservices với ba dịch vụ độc lập đòi hỏi sự giao tiếp và hợp tác chặt chẽ giữa các nhóm, từ khâu thống nhất API Contract, xử lý xác thực JWT cho đến đồng bộ dữ liệu qua cơ chế Event. Tinh thần làm việc nhóm và sự hỗ trợ lẫn nhau giữa ba nhóm là yếu tố then chốt dẫn đến thành công của toàn bộ dự án. Cuối cùng, nhóm xin gửi lời cảm ơn đến gia đình và bạn bè đã luôn động viên, ủng hộ và tạo điều kiện tốt nhất để mỗi thành viên có thể dành trọn thời gian và tâm huyết cho bài tập lớn này. Mặc dù đã nỗ lực hết mình, song do hạn chế về thời gian và kinh nghiệm, báo cáo không tránh khỏi những thiếu sót. Nhóm em rất mong nhận được những ý kiến đóng góp quý báu từ thầy và các bạn để có thể hoàn thiện hơn nữa.

Nhóm N2 – Circulation Service Hà Nội, tháng 7 năm 2026
